\documentclass[10pt]{article}
\usepackage[letterpaper,margin=.65in]{geometry}
\usepackage[T1]{fontenc}
\usepackage{lmodern,amsmath,graphicx,booktabs,threeparttable,float,adjustbox,etoolbox}
\usepackage[colorlinks=true,linkcolor=blue,urlcolor=blue]{hyperref}
\usepackage{caption}
\captionsetup{font=small,labelfont=bf}
\renewcommand{\TPTminimum}{\linewidth}
\setlength{\parindent}{0pt}
\setlength{\parskip}{5pt}
\BeforeBeginEnvironment{tabular}{\begin{adjustbox}{max width=\linewidth}}
\AfterEndEnvironment{tabular}{\end{adjustbox}}
\let\originaltable\table
\renewcommand{\table}[1][]{\originaltable[H]\small}
\newcommand{\entry}[3]{\clearpage\section{#1}\textbf{Source:} \path{#2}\par #3\par}
\begin{document}
\begin{center}
{\LARGE Young Lives: local results viewer}\\[6pt]
{\large Personal research notes --- not for publication}\\[4pt]
Compiled \today
\end{center}
This project reads current staged tables from \path{../output/tables/}. It does not rerun regressions when compiled. The build script refreshes the employment figure using the original plotting expression and existing processed data. Outputs may originate from different analysis runs; compilation date is not an estimation date. Source headers record table timestamps where available.

\textbf{How to refresh.} From the repository root run \path{./results_viewer/build.sh}. To regenerate all analysis outputs first, run \path{Rscript scripts/master.R}; this overwrites processed data and results, including manuscript Table 3. Table 2 alone: \path{Rscript scripts/table_main.R}. Then rebuild this viewer. Open \path{results_viewer/main.tex} in a LaTeX editor or read \path{results_viewer/main.pdf}.

\textbf{Reading guide.} Viewer table numbers differ from manuscript numbers. Section titles identify manuscript equivalents. All tables show the files as generated; the notes identify known export issues rather than silently changing their headings or units. Table 2 is reconstructed and verified; the legacy exports for manuscript Tables 4 and 5 are not exact reproductions of the December 2023 manuscript versions.

\textbf{Common design.} Older-cohort data are built by \path{scripts/extract_clean.R}. Models generally filter \path{raw_lang != 0}, then use complete cases for each outcome and its covariates. Main 2SLS models use \path{lfe::felm}: $y=X\gamma+\beta IMTI+u$, with $IMTI$ instrumented by $E_{is}$. The salary variable is \path{wage_employII}; broader wage employment is \path{wage_employ} and includes salary positives. Missingness and exclusions differ. Conventional model standard errors are not region-clustered unless explicitly stated. Table 2 now includes CR2/Satterthwaite inference and a cluster sign-flip AR sensitivity check. Few-region inference remains fragile.

\textbf{Main IV controls (X).} Child sex, BMI-for-age z-score, stunting, caregiver education, age and sex, household size, wealth, housing quality, consumer durables, electricity, land ownership and food-security categories. The main IV models have no region fixed effects. Standardized scores are constructed separately for Addis and non-Addis samples before outcome-specific missing-value exclusions.
\tableofcontents

\entry{Summary statistics --- manuscript Table 1}{scripts/descriptives.R}{No regression: means, available-case counts and standard deviations, separately for non-Addis and Addis samples. Region indicators enter the non-Addis summary only. Input: the two older-cohort CSVs. Export: \path{output/tables/tableI.tex}.}
% latex table generated in R 4.6.0 by xtable 1.8-8 package
% Sat Oct 31 16:41:05 2026
\begin{table}[ht]
\centering
\caption{Summary Statistics} 
\label{tab:sum-stat}
\begin{tabular}{llrrrlrrr}
  \toprule
   \\[-1.8ex]  & & \multicolumn{3}{c}{Non AA Sample} & & \multicolumn{3}{c}{AA Sample}  \\[0.2ex] \cline{3-5}  \cline{7-9}  \\[-1.2ex]  & & \multicolumn{1}{c}{N} & \multicolumn{1}{c}{Mean} & \multicolumn{1}{c}{sd} & & \multicolumn{1}{c}{N} & \multicolumn{1}{c}{Mean} & \multicolumn{1}{c}{sd} \\  \\[-1.8ex] &  & \multicolumn{1}{c}{(1)} & \multicolumn{1}{c}{(2)} & \multicolumn{1}{c}{(3)} & & 
  \multicolumn{1}{c}{(4)} & \multicolumn{1}{c}{(5)} & \multicolumn{1}{c}{(6)} \\  \midrule
 & Dummy for male child &   728 & 0.544 &  &  &   130 & 0.492 &  \\ 
   & Highest grade completed &   723 & 9.264 & 3.143 &  &   130 & 11.26 & 2.813 \\ 
   & Maths test Score (Raw) &   725 & 13.12 &  5.77 &  &   130 & 17.03 & 4.905 \\ 
   & Language test Score (Raw) &   728 & 13.88 & 4.306 &  &   130 & 17.48 & 3.344 \\ 
   & Dummy for wage Employment &   310 & 0.497 & 0.501 &  &    86 & 0.605 & 0.492 \\ 
   & Dummy for salaried Employment &   487 & 0.162 & 0.369 &  &    88 & 0.386 &  0.49 \\ 
   & BMI-for-age z-score &   673 & -1.682 &  1.82 &  &   128 & -0.971 & 1.407 \\ 
   & Dummy for stunting &   700 & 0.213 &  &  &   128 &  0.18 &  \\ 
   & Dummy for female caregiver &   727 & 0.978 &  &  &   130 & 0.985 &  \\ 
   & Caregiver's age &   727 & 35.17 & 8.998 &  &   130 & 36.59 & 9.709 \\ 
   & Caregiver's highest grade &   719 &  2.73 & 5.899 &  &   129 & 4.403 & 5.269 \\ 
   & Household size &   727 & 6.399 & 2.082 &  &   129 &  6.07 & 2.312 \\ 
   & Wealth index &   724 & 0.343 & 0.168 &  &   129 & 0.576 & 0.107 \\ 
   & Housing quality index &   726 & 0.327 & 0.195 &  &   129 & 0.417 &  0.14 \\ 
   & Consumer durables index &   727 & 0.257 & 0.197 &  &   129 & 0.492 & 0.157 \\ 
   & Household has access to electricity &   726 & 0.482 &   0.5 &  &   129 & 0.977 & 0.151 \\ 
   & Household owns land where house is on &   728 & 0.779 & 0.415 &  &   130 & 0.154 & 0.362 \\ 
   & Household resides in a rural area &   727 & 0.726 &  &  &  &  &  \\ 
   & Dummy for Amhara &   728 & 0.236 &  &  &  &  &  \\ 
   & Dummy for Oromia &   728 & 0.238 &  &  &  &  &  \\ 
   & Dummy for SNNP &   728 & 0.268 &  &  &  &  &  \\ 
   & Dummy for Tigray &   728 & 0.258 &  &  &  &  &  \\ 
    \bottomrule   
 \\[-1.8ex] \multicolumn{6}{l}{\footnotesize{\textit{Note:} The standard deviations for proportions is 
               not presented.}} 
\end{tabular}
\end{table}


\entry{Employment and instrument diagnostics --- manuscript Table 2}{scripts/table_main.R}{Four 2SLS regressions: salary and wage outcomes, whole non-Addis sample and excluding Amhara; common IV controls listed above. Each diagnostic uses its exact second-stage sample. Standalone generation or automatic call from \path{scripts/export_res.R}. Full-precision diagnostics: \path{output/tables/tableII_mod_diagnostics.csv}. Panel C reports regional CR2/Satterthwaite inference and a weak-IV-robust sign-flip sensitivity check; all four CR2 confidence intervals include zero. The original conventional stars in Panel A are retained only for comparison. The salary significance is not robust to clustering. See the clustering notes below.}
% Generated by scripts/table_main.R; no manual numerical edits.
\begin{table}[!htbp] \centering
\linespread{1}\selectfont
\caption{Main results} \label{tab:mainreg}
\begin{threeparttable}
\small
\begin{tabular}{@{\extracolsep{5pt}}lcccc}
\hline\hline
 & \multicolumn{2}{c}{\textbf{Whole Sample}} & \multicolumn{2}{c}{\textbf{Excluding Amhara}} \\
\cline{2-3} \cline{4-5}
 & Salary & Wage & Salary & Wage \\
 & Employed & Employed & Employed & Employed \\
 & (1) & (2) & (3) & (4) \\
\hline
\multicolumn{5}{l}{\textbf{Panel A: 2SLS (conventional inference)}} \\
$MTIuse$ & 0.043$^{***}$ & 0.029 & 0.042$^{**}$ & 0.029 \\
 & (0.016) & (0.027) & (0.018) & (0.029) \\
\hline
\multicolumn{5}{l}{\textbf{Panel B: First stage}} \\
$E_{is}$ & 3.721$^{***}$ & 3.641$^{***}$ & 3.627$^{***}$ & 3.404$^{***}$ \\
 & (0.297) & (0.388) & (0.333) & (0.463) \\
Excluded-instrument $F$ & 157.07 & 87.97 & 118.79 & 54.17 \\
Partial $R^2$ & 0.272 & 0.251 & 0.273 & 0.236 \\
First-stage $F$ (HC1) & 81.48 & 41.91 & 67.55 & 28.80 \\
\hline
\multicolumn{5}{l}{\textbf{Panel C: Regional clustering (small-sample checks)}} \\
2SLS SE (CR2) & 0.018 & 0.009 & 0.017 & 0.017 \\
2SLS $p$ (CR2/Satt.) & 0.172 & 0.094 & 0.194 & 0.276 \\
2SLS Satterthwaite d.f. & 1.62 & 1.95 & 1.31 & 1.44 \\
AR $p$ (CR2/Satt.) & 0.123 & 0.274 & 0.081 & 0.438 \\
AR $p$ (cluster sign-flip) & 0.250 & 0.250 & 0.500 & 0.500 \\
Informative regions (sign-flip) & 3 & 3 & 2 & 2 \\
Regions / sign patterns & 4 / 16 & 4 / 16 & 3 / 8 & 3 / 8 \\
\hline
Observations & 437 & 279 & 334 & 192 \\
\hline\hline
\end{tabular}
\begin{tablenotes} \footnotesize
\item \textit{Notes:} Controls are child sex, BMI-for-age z-score, stunting, caregiver education, age and sex, household size, wealth, housing-quality and consumer-durables indices, electricity, land ownership, and food-security categories. Each first stage uses the exact corresponding second-stage sample.
\item Panels A--B retain conventional/HC1 inference for comparison. First-stage $F$ tests the excluded instrument conditional on controls. Partial $R^2$ measures its incremental explanatory power. One instrument identifies one endogenous regressor; overidentification cannot be tested.
\item Panel C: region-clustered CR2 with coefficient-specific Satterthwaite d.f.; Wald inference also requires adequate identification. AR tests $H_0:\beta_{MTI}=0$. Sign-flip AR residualizes the instrument on controls within each region and enumerates all Rademacher patterns. It discards between-region variation; Tigray has no residual instrument variation. With only three/two informative regions, minimum two-sided $p$-values are .250/.500. Independent regions, conditional instrument validity and approximately symmetric cluster scores are required; enumeration is not an exact finite-sample guarantee. No Panel C test rejects at 5\%; salary significance is not robust to clustering.
\item Conventional standard errors are in parentheses. *** Significant at 1\%, ** at 5\%, * at 10\%.
\end{tablenotes}
\end{threeparttable}
\end{table}


\entry{Cognitive channel --- manuscript Table 3}{scripts/table_rev.R}{OLS via \path{felm}, not a completed 2SLS channel analysis: $std\_lang=X\gamma+\delta(E_{is}\times IMTI)+u$. Same IV controls; non-Addis sample, whole and excluding Amhara. Language scores are standardized before excluding Amhara. This input comes from the local output folder; the generator also mirrors it to Dropbox.}
\input{../output/tables/table_rev.tex}

\entry{Addis Ababa --- manuscript Table 4 counterpart}{scripts/estimation.R; scripts/export_res.R}{Current staged export, not the historical manuscript table. Panel A uses OLS reduced forms ($y$ on $E_{is}$); Panel B uses 2SLS. Sample excludes activity code 19 and filters zero language scores. IV controls add school type (\path{entype_r4}). Reduced-form controls differ: BMI, caregiver characteristics, household indices, sex, child relationship, preschool, age at grade 1, food security and school type. Exported model order: \path{rf6aa, rf7aa, rf5aa, rf1aa} and \path{iv8aa, iv9aa, iv6aa, iv1aa}.}

%Tex File: /Users/user/Library/CloudStorage/GoogleDrive-tayelemi@gmail.com/My Drive/Personal/R_projects/potential_research-Young_Lives/output/tables/tableIII.tex

% Date and time: Sat, Oct 31, 2026 - 16:41:07
\begin{table}[!htbp] \centering 
  \caption{Estimates for Addis Ababa} 
  \label{tab:aareg} 
\begin{threeparttable}
\begin{tabular}{@{\extracolsep{5pt}}lcccc} 
\\[-1.8ex]\hline 
\hline \\[-1.8ex] 
 & Maths & Language & Salary & Wage \\ 
 & Z-score & Z-score & Employed & Employed \\ 
\\[-1.8ex] & (1) & (2) & (3) & (4)\\ 
\hline \\[-1.8ex] 
\\[-2.0ex] \multicolumn{5}{@{} l}{\textbf{Panel A: Reduced Form}}
 \\
 \\[-1.5ex]
 $E_{is}$ & 0.667 & 0.025 & $-$0.313 & 0.000 \\ 
  & (0.492) & (0.411) & (0.200) & (0.000) \\ 
  & & & & \\ 
\\[-1.83ex] 
 \hline \\[-1.83ex]
\\[-2.0ex] \multicolumn{5}{@{} l}{\textbf{Panel B: 2SLS}}
 \\
 \\[-1.5ex]
 IMTI & $-$0.041 & $-$0.045 & 0.0002 & 0.000 \\ 
  & (0.067) & (0.052) & (0.031) & (0.000) \\ 
  & & & & \\ 
\hline \\[-1.8ex] 
Observations & 49 & 49 & 41 & 39 \\ 
\hline 
\hline \\[-1.8ex] 
\end{tabular} 
\begin{tablenotes}
\footnotesize
\item \textit{Notes:} Standard errors are in parentheses. *** Significant at 1\%, ** Significant at 5\%, * Significant at 10\%.
\end{tablenotes}
\end{threeparttable}
\end{table} 





\entry{Tigray versus Oromia --- manuscript Table 5 counterpart}{scripts/estimation.R; scripts/export_res.R}{OLS regressions on a Tigray indicator, controlling for IMTI, land ownership, rural residence and the first-stage covariate set (BMI, caregiver characteristics, household indices, sex, child relationship, preschool, age at grade 1 and food security). Verbal regressions additionally require test language to match the child's language. Models: \path{rf6ot, rf7ot, rf5ot, rf1ot}. \textbf{Known heading issue:} columns 3 and 4 are actually salary then wage, although this export labels them Wage then Salary. This older export also differs from the manuscript version.}

%Tex File: /Users/user/Library/CloudStorage/GoogleDrive-tayelemi@gmail.com/My Drive/Personal/R_projects/potential_research-Young_Lives/output/tables/tableV.tex

% Date and time: Sat, Oct 31, 2026 - 16:41:07
\begin{table}[!htbp] \centering 
  \caption{Comparing Estimate for Tigray and Oromia} 
  \label{tab:oromtig} 
\begin{threeparttable}
\begin{tabular}{@{\extracolsep{5pt}}lcccc} 
\\[-1.8ex]\hline 
\hline \\[-1.8ex] 
\\[-1.8ex] & Maths & Language & Wage & Salary \\ 
 & Z-score & Z-score & Employed & Employed \\
\\[-1.8ex] & (1) & (2) & (3) & (4)\\ 
\hline \\[-1.8ex] 
 Tigray & $-$0.072 & 0.210 & 0.113 & 0.086 \\ 
  & (0.191) & (0.164) & (0.099) & (0.147) \\ 
  & & & & \\ 
\hline \\[-1.8ex] 
Observations & 331 & 283 & 227 & 130 \\ 
R$^{2}$ & 0.255 & 0.446 & 0.127 & 0.276 \\ 
\hline 
\hline \\[-1.8ex] 
\end{tabular} 
\begin{tablenotes}
\footnotesize
\item \textit{Notes:} Standard errors are in parentheses. *** Significant at 1\%, ** Significant at 5\%, * Significant at 10\%.
\end{tablenotes}
\end{threeparttable}
\end{table} 





\entry{Legacy non-Addis export --- supplementary}{scripts/estimation.R; scripts/export_res.R}{First-stage and 2SLS panels for standardized math, standardized verbal, salary and wage outcomes. Main IV controls. Models: \path{iv8, iv9, iv6, iv1}; Panel A draws their \path{stage1} models. Superseded for employment reporting by the updated manuscript Table 2, but retained here to inspect generated outputs. The original and updated tables reuse a label; this viewer suppresses the duplicate below.}
\begingroup\renewcommand{\label}[1]{}
%Tex File: /Users/user/Library/CloudStorage/GoogleDrive-tayelemi@gmail.com/My Drive/Personal/R_projects/potential_research-Young_Lives/output/tables/tableII.tex

% Date and time: Sat, Oct 31, 2026 - 16:41:07
\begin{table}[!htbp] \centering 
  \caption{Main Results} 
  \label{tab:mainreg} 
\begin{threeparttable}
\begin{tabular}{@{\extracolsep{5pt}}lcccc} 
\\[-1.8ex]\hline 
\hline \\[-1.8ex] 
 & Maths & Language & Salary & Wage \\ 
 & Z-score & Z-score & Employed & Employed \\
\\[-1.8ex] & (1) & (2) & (3) & (4)\\ 
\hline \\[-1.8ex] 
\\[-2.0ex] \multicolumn{5}{@{} l}{\textbf{Panel A: First Stage}}
 \\
 \\[-1.5ex]
 $E_{is}$ & 3.684$^{***}$ & 3.665$^{***}$ & 3.721$^{***}$ & 3.641$^{***}$ \\ 
  & (0.257) & (0.257) & (0.297) & (0.388) \\ 
  & & & & \\ 
\\[-1.83ex] 
 \hline \\[-1.83ex]
\\[-2.0ex] \multicolumn{5}{@{} l}{\textbf{Panel B: 2SLS}}
 \\
 \\[-1.5ex]
 IMTI & 0.058$^{*}$ & 0.066$^{**}$ & 0.043$^{***}$ & 0.029 \\ 
  & (0.030) & (0.030) & (0.016) & (0.027) \\ 
  & & & & \\ 
\hline \\[-1.8ex] 
Observations & 654 & 656 & 437 & 279 \\ 
\hline 
\hline \\[-1.8ex] 
\end{tabular} 
\begin{tablenotes}
\footnotesize
\item \textit{Notes:} Standard errors are in parentheses. *** Significant at 1\%, ** Significant at 5\%, * Significant at 10\%.
\end{tablenotes}
\end{threeparttable}
\end{table} 



\endgroup

\entry{Younger cohort, Addis Ababa --- supplementary}{scripts/cleaning_yc.R; scripts/estimation.R; scripts/export_res.R}{Uses \path{ylsample_yc.csv} and \path{schsur_yl_joined.csv}. First two columns: Young Lives math/verbal; last two: school-survey math/literacy. Young Lives models restrict standardized math below 2.25. Panel A uses standardized outcomes; Panel B uses raw test scores despite the repeated ``Z-score'' export heading. Controls include school type, sex, BMI, stunting, caregiver education/age, household size and indices, electricity, land and food security. Caregiver sex is not included here. Do not interpret differences across panels as a common standardized effect.}

%Tex File: /Users/user/Library/CloudStorage/GoogleDrive-tayelemi@gmail.com/My Drive/Personal/R_projects/potential_research-Young_Lives/output/tables/tableIV.tex

% Date and time: Sat, Oct 31, 2026 - 10:30:12
\begin{table}[!htbp] \centering 
  \caption{Estimates for Addis Ababa: Younger Cohort} 
  \label{tab:ycaa} 
\begin{threeparttable}
\begin{tabular}{@{\extracolsep{5pt}}lcccc} 
\\[-1.8ex]\hline 
\hline \\[-1.8ex] 
 & \multicolumn{2}{c}{\textbf{YL Cognitive}} & \multicolumn{2}{c}{\textbf{School Survey}} \\ 
\cline{2-3} \cline{4-5} 
\\[-1.8ex]
\\[-1.8ex] & Maths & Language & Maths & Literacy \\ 
 & Z-score & Z-score & Z-score & Z-score \\
\\[-1.8ex] & (1) & (2) & (3) & (4)\\ 
\hline \\[-1.8ex] 
\\[-2.0ex] \multicolumn{5}{@{} l}{\textbf{Panel A: Reduced Form}}
 \\
 \\[-1.5ex]
 $E_{is}$ & 0.209 & 0.147 & $-$0.218 & $-$0.107 \\ 
  & (0.179) & (0.192) & (0.421) & (0.411) \\ 
  & & & & \\ 
\\[-1.83ex] 
 \hline \\[-1.83ex]
\\[-2.0ex] \multicolumn{5}{@{} l}{\textbf{Panel B: 2SLS}}
 \\
 \\[-1.5ex]
 IMTI & 0.133 & 0.796 & $-$0.169 & $-$0.095 \\ 
  & (0.114) & (1.036) & (0.327) & (0.368) \\ 
  & & & & \\ 
\hline \\[-1.8ex] 
Observations & 249 & 247 & 71 & 71 \\ 
\hline 
\hline \\[-1.8ex] 
\end{tabular} 
\begin{tablenotes}
\footnotesize
\item \textit{Notes:} Standard errors are in parentheses. *** Significant at 1\%, ** Significant at 5\%, * Significant at 10\%.
\end{tablenotes}
\end{threeparttable}
\end{table} 





\entry{Excluding Amhara --- supplementary}{scripts/estimation.R; scripts/export_res.R}{Same first-stage and 2SLS design as the legacy non-Addis export, restricted to regions other than Amhara. Main IV controls. Models: \path{iv8_nonAm, iv9_nonAm, iv6_nonAm, iv1_nonAm}. Employment columns are reproduced in the revised Table 2.}

%Tex File: /Users/user/Library/CloudStorage/GoogleDrive-tayelemi@gmail.com/My Drive/Personal/R_projects/potential_research-Young_Lives/output/tables/tableVI.tex

% Date and time: Sat, Oct 31, 2026 - 16:41:07
\begin{table}[!htbp] \centering 
  \caption{Estimates Excluding Observations in the Amhara Region} 
  \label{tab:nonAm} 
\begin{threeparttable}
\begin{tabular}{@{\extracolsep{5pt}}lcccc} 
\\[-1.8ex]\hline 
\hline \\[-1.8ex] 
 & Maths & Language & Salary & Wage \\ 
 & Z-score & Z-score & Employed & Employed \\
\\[-1.8ex] & (1) & (2) & (3) & (4)\\ 
\hline \\[-1.8ex] 
\\[-2.0ex] \multicolumn{5}{@{} l}{\textbf{Panel A: First Stage}}
 \\
 \\[-1.5ex]
 $E_{is}$ & 3.568$^{***}$ & 3.551$^{***}$ & 3.627$^{***}$ & 3.404$^{***}$ \\ 
  & (0.284) & (0.285) & (0.333) & (0.463) \\ 
  & & & & \\ 
\\[-1.83ex] 
 \hline \\[-1.83ex]
\\[-2.0ex] \multicolumn{5}{@{} l}{\textbf{Panel B: 2SLS}}
 \\
 \\[-1.5ex]
 IMTI & 0.061$^{*}$ & 0.061$^{*}$ & 0.042$^{**}$ & 0.029 \\ 
  & (0.032) & (0.033) & (0.018) & (0.029) \\ 
  & & & & \\ 
\hline \\[-1.8ex] 
Observations & 522 & 524 & 334 & 192 \\ 
\hline 
\hline \\[-1.8ex] 
\end{tabular} 
\begin{tablenotes}
\footnotesize
\item \textit{Notes:} Standard errors are in parentheses. *** Significant at 1\%, ** Significant at 5\%, * Significant at 10\%.
\end{tablenotes}
\end{threeparttable}
\end{table} 





\entry{Employment by region --- manuscript Figure 1}{scripts/descriptives.R; results_viewer/refresh_figure.R}{No regression. Available-case means of salary and broader wage indicators in the filtered non-Addis sample. The viewer evaluates the original \path{employ_gr} expression directly from \path{descriptives.R}, then saves it with base R's PDF device to avoid the local Cairo/XQuartz dependency. Bar values are unchanged; fonts may differ. Confidence intervals and bar-specific sample sizes have not yet been added.}
\begin{figure}[H]\centering
\includegraphics[width=.95\linewidth]{assets/employ_gr.pdf}
\caption{Percentage employed by Round 5, by region}
\end{figure}
\clearpage\section{Regional clustering: methods and conclusions}
\textbf{Reproduce:} \path{Rscript scripts/table_main.R}. Helper:
\path{scripts/cluster_inference.R}. Checks:
\path{Rscript tests/check_table2_clustering.R}.
Numerical results: \path{output/tables/tableII_mod_diagnostics.csv}.
Region-level sample sizes and AR scores: \path{output/tables/tableII_mod_cluster_audit.csv}.

All original 2SLS point estimates, samples and controls are retained.
Clustering changes the inference. CR2 bias-adjusted regional variances use
coefficient-specific Satterthwaite degrees of freedom, implemented by
\texttt{clubSandwich} on equivalent \texttt{AER::ivreg} fits. Each fit is checked
against the original \texttt{felm} estimate. These are small-sample approximations,
not exact tests, and Wald inference also requires adequate identification.

\begin{center}
\begin{tabular}{lrrrr}
\toprule
 & \multicolumn{2}{c}{Whole sample} & \multicolumn{2}{c}{Excluding Amhara}\\
 & Salary & Wage & Salary & Wage\\\midrule
Estimate & .0431 & .0292 & .0423 & .0289\\
Conventional $p$ & .007 & .277 & .018 & .322\\
CR2 SE & .0183 & .0095 & .0171 & .0169\\
Satterthwaite d.f. & 1.62 & 1.95 & 1.31 & 1.44\\
CR2 $p$ & .172 & .094 & .194 & .276\\
AR CR2 $p$ & .123 & .274 & .081 & .438\\
AR sign-flip $p$ & .250 & .250 & .500 & .500\\\bottomrule
\end{tabular}
\end{center}
\textbf{Conclusion.} Salary employment no longer rejects a zero effect at 5\%.
None of the four CR2 tests rejects at 5\%; all CR2 95\% confidence intervals include
zero. Whole-sample wage is marginal at 10\% under the Wald check, while the
corresponding AR test is not. Conversely, the excluding-Amhara salary AR test is
marginal at 10\%, while its Wald test is not. Neither is stable across methods.
The evidence cannot establish a robust positive employment effect or a precise zero.

\textbf{Weak-instrument check.} The CR2 AR test applies the same adjustment to
$y$ on controls and $E_{is}$ under $H_0:\beta=0$. It avoids dividing by an estimated
first-stage coefficient, but still relies on small-sample approximations.
Clustered first-stage F statistics, d.f. and p-values are recorded in the CSV;
with 3--4 regions they should not be assessed using many-cluster weak-IV cutoffs.

\textbf{Sign-flip sensitivity.} Inspired by Wang--Zhang's wild AR score procedure,
the instrument is orthogonalized against controls separately by region to avoid
assuming the same instrument--control projection across regions. Null-restricted
scores are aggregated by region and all $2^G$ Rademacher patterns are enumerated.
This is a different set of test moments from the pooled instrument: it discards
between-region variation, while leaving the reported point estimator unchanged.
Tigray's residual instrument is zero; only three regions, or two excluding Amhara,
are informative. Tied patterns are included. Attainable minimum two-sided
p-values are .250 and .500. Additional draws would merely repeat these patterns.

The check requires independent regions, conditional instrument validity within
regions, and approximately symmetric cluster scores, supported asymptotically
by suitable within-region weak dependence. It does not accommodate arbitrary
common regional factor shocks, justify the exclusion restriction, or provide an
exact finite-sample guarantee. Policy assignments are not permuted: no known
random assignment mechanism is assumed. Webb weights or smaller administrative
clusters should not be selected simply to obtain finer or smaller p-values.

\textbf{Method sources:}
\href{https://doi.org/10.1080/07350015.2016.1247004}{Pustejovsky and Tipton (2018)},
\href{https://jepusto.github.io/clubSandwich/}{clubSandwich documentation}, and
\href{https://doi.org/10.1016/j.jeconom.2024.105727}{Wang and Zhang (2024)}.

\end{document}
